\documentclass[11pt]{article}
\input{packages.tex}
\usepackage{diagbox} % Пакет для диагонального разделения ячейки
\usepackage{tikz}
\DeclareMathOperator{\sign}{sign}
\DeclareMathOperator{\Imm}{Im}

% Настройка колонтитулов
\pagestyle{fancy}
\fancyhf{} % Очистить все поля
\fancyhead[L]{Машины, автоматы, языки... Картина маслом!}
\fancyhead[R]{\href{https://t.me/ceagest}{Автор}}
\fancyfoot[C]{\thepage}               % Справа внизу

\setlength{\headheight}{13.6pt}

\hypersetup{
    colorlinks=true,
    linkcolor=blue,
    citecolor=blue,
    filecolor=blue,
    urlcolor=blue,
}

\newcommand{\rdot}[1]{%
    \tikz[remember picture, baseline=-0.5ex] \node[circle, fill=red, inner sep=1.5pt] (#1) {};%
}

\begin{document}

\begin{titlepage}
    \centering
    \vspace*{\fill}  % заполняет пространство сверху

    {\Huge Лекция 11. Разрешимость и перечислимость.
    \par}

    \vspace*{\fill}  % заполняет пространство снизу
\end{titlepage}

\setcounter{page}{1}

\section{Разрешимость и полуразрешимость}

\begin{definition}
    Язык $\mathcal{L}$ в конечном алфавите $A$ называется \textbf{разрешимым}, если существует такая решающая МТ с алфавитом $B \supset A$, которая:
    \begin{itemize}
        \item на любом входе $w \in \mathcal{L}$ останавливается в состоянии $q_{\text{yes}}$
        \item на любом входе $w \not \in \mathcal{L}$ останавливается в состоянии $q_{\text{no}}$
    \end{itemize}
    При этом пустой символ $\Lambda$ не принадлежит алфавиту $A$. Напомним, что \textbf{решающая} машина Тьюринга имеет ровно два финальных состояния: $q_{\text{yes}}$ и $q_{\text{no}}$.
\end{definition}

\begin{definition}
    Язык $\mathcal{L}$ в конечном алфавите $A$ называется \textbf{допустимым (полуразрешимым)}, если существует такая решающая МТ с алфавитом $B \supset A$, которая:
    \begin{itemize}
        \item на любом входе $w \in \mathcal{L}$ останавливается в состоянии $q_{\text{yes}}$
        \item на любом входе $w \not \in \mathcal{L}$ \textbf{НЕ} останавливается в состоянии $q_{\text{yes}}$ (для удобства считаем, что не останавливается вовсе)
    \end{itemize}
    При этом пустой символ $\Lambda$ не принадлежит алфавиту $A$.
\end{definition}

\begin{lemma}
    Существуют полуразрешимые языки, не являющиеся разрешимыми.
\end{lemma}

\begin{proof}
    В самом деле, рассмотрим язык \textbf{stop} --- язык в алфавите $\{0, 1\}$, состоящий из тех описаний $\left\langle M, w \right\rangle$ машины Тьюринга $M$ и входного слова $w$, для которых машина $M$ останавливается на входе $w$. Из проблемы остановки следует, что язык \textbf{stop} не является разрешимым. При этом \textbf{stop} является полуразрешимым, так как в качестве полурешающей МТ для этого языка можно взять универсальную машину Тьюринга: для заданного слова $\left\langle M, w \right\rangle \in \text{\textbf{stop}}$ универсальная МТ моделирует работу $M$ для входа $w$, останавливаясь в состоянии $q_{\text{yes}}$, если машина $M$ остановилась на входе $w$, и не останавливаясь вовсе, если машина $M$ не остановилась на входе $w$. Таким образом, лемма доказана.
\end{proof}

\section{Полуразрешимость и перечислимость}

\begin{definition}
    Язык $\mathcal{L}$ называется \textbf{перечислимым}, если существует МТ, которая, работая с пустого входа, напечатает на выходной ленте все слова языка $\mathcal{L}$, разделяя их специальными символами. При этом перечисляющая МТ никогда не изменяет непустые символы выходной ленты.
\end{definition}

\begin{remark}
    Язык $\mathcal{L}$ чаще всего является бесконечным. Под "напечатает на выходной ленте все слова языка $\mathcal{L}$" \, подразумевается, что каждое слово данного языка появится на ленте за \textbf{конечное} время. При этом заметим, что нет никаких гарантий на порядок, в котором будут печататься слова. Более того, не требуется даже отсутствия повторов слов (но, при соответствующем желании, уникальность записи слов можно реализовать, что будет сделано позже).
\end{remark}

\begin{remark}
    Под \textbf{входной} лентой МТ будем понимать ленту, с которой возможно лишь читать данные. Под \textbf{выходной} лентой МТ понимается лента, на которой возможно лишь один раз записать данные слева направо. Как входная, так и выходная ленты МТ \textbf{НЕ} учитываются как память этой МТ. Память учитывается только на \textbf{рабочих} лентах МТ --- лентах, которые не совпадают со входной и выходной лентами. Если машина Тьюринга имеет лишь одну ленту, то она априорно полагается рабочей, даже если на неё записывается вход. 
\end{remark}

\begin{theorem}
    Язык полуразрешим тогда и только тогда, когда он перечислим.
\end{theorem}

\begin{proof}
    $(\Longrightarrow):$ Пусть язык $\mathcal{L}$ над алфавитом $\Sigma$ полуразрешим. Тогда для него существует полурешающая машина Тьюринга $M_1$. Построим на основе $M_1$ машину Тьюринга $M_2$, перечисляющую $\mathcal{L}$. Она будет иметь пять лент: на первой перебираются всевозможные слова языка над алфавитом $\Sigma$, на второй эмулируется работа $M_1$ на соответствующих входах, третья используется для подсчитывания количества операций (шагов), совершаемых при эмуляции работы $M_1$ (это сделано для того, чтобы сравнивать требуемое число шагов в ячейке таблицы с текущим совершённым в процессе эмуляции их количеством), на четвёртой поддерживаются счётчики обхода змейкой (это нужно для того, чтобы понимать, куда змейка идёт сейчас --- направо, налево, вверх или вниз), а на пятой печатаются те слова, которые приняла $M_1$ (то есть требуемое перечисление). Перебирать слова и запускать на них $M_1$ мы будем специальным образом: составим бесконечную таблицу, в первом столбце которой будут записаны все слова над $\Sigma$ (коих, ясен пень, счётно), идущие в лексикографическом порядке, а в первой строке --- количество операций, которое (или не более, чем которое) машина $M_1$ должна совершить на соответствующем входном слове. Проходить эту таблицу мы будем уже известным всем нам из матанализа способом, а именно змейкой, начинающей свой путь в левой верхней клетке таблицы, а далее идущей так, как показано на рисунке. \\
    \begin{table}[htbp]
    \centering
    \renewcommand{\arraystretch}{2}
    \begin{tabular}{|c|c|c|c|c|c|}
        \hline
        \diagbox{Слово}{Кол-во операций} & 1000 & 2000 & 3000 & 4000 & $\dots$ \\
        \hline
        $w_1$ & \rdot{n11} & \rdot{n12} & \rdot{n13} & \rdot{n14} & \\
        \hline
        $w_2$ & \rdot{n21} & \rdot{n22} & \rdot{n23} & \rdot{n24} & \\
        \hline
        $w_3$ & \rdot{n31} & \rdot{n32} & \rdot{n33} & \rdot{n34} & \\
        \hline
        $w_4$ & \rdot{n41} & \rdot{n42} & \rdot{n43} & \rdot{n44} & \\
        \hline
        $\dots$ & & & & & \\
        \hline
    \end{tabular}
    \label{tab:snake_time}
    \end{table}

    \begin{tikzpicture}[remember picture, overlay, >=stealth]
    \tikzset{snake path/.style={red, thick, ->}}

    \draw[snake path] (n11) -- (n21);
    \draw[snake path] (n21) -- (n22);
    \draw[snake path] (n22) -- (n12);

    \draw[snake path] (n12) -- (n13);
    \draw[snake path] (n13) -- (n23);
    \draw[snake path] (n23) -- (n33);
    \draw[snake path] (n33) -- (n32);
    \draw[snake path] (n32) -- (n31);

    \draw[snake path] (n31) -- (n41);
    \draw[snake path] (n41) -- (n42);
    \draw[snake path] (n42) -- (n43);
    \draw[snake path] (n43) -- (n44);
    \draw[snake path] (n44) -- (n34);
    \draw[snake path] (n34) -- (n24);
    \draw[snake path] (n24) -- (n14);

    \draw[red, thick, dashed, ->] (n14) -- ++(1, 0); 
    \end{tikzpicture} \leavevmode
    \\
    Ясно, что таким образом посещены будут все клетки таблицы, а значит для каждого слова над $\Sigma$ будет запущена $M_1$ на всех возможных ограничениях количества выполняемых ею на этом входе шагов и, следовательно, если рассматриваемое слово присутствует в языке, то за конечное число шагов змейка попадёт в клетку с таким числом операций, для которого слово будет принято $M_1$ и напечатано в перечислении. Если же рассматриваемого слова в языке нет, то при любом ограничении на количество шагов машина $M_1$ за это число операций не остановится в $q_{\text{yes}}$, а значит и в перечислении этого слова никогда не появится. \\
    $(\Longleftarrow):$ Пусть язык $\mathcal{L}$ перечислим. Тогда существует машина Тьюринга $M$, его перечисляющая. Построим на основе машины $M$ новую машину Тьюринга $M'$, которая будет являться полурешающей для языка $\mathcal{L}$. Пусть $M'$ на входе $x$ эмулирует работу $M$ с пустого входа (так как перечисляющая МТ с него только и работает) и сравнивает выводимые $M$ слова со словом $x$. Если произошло совпадение, машина $M'$ останавливается в $q_{\text{yes}}$, иначе продолжает работу. Таким образом, если рассматриваемое слово лежит в языке, то оно по определению есть и в перечислении, а значит $M'$ остановится на нём в состоянии $q_{\text{yes}}$. Если же рассматриваемого слова в языке нет, то его и в перечислении не встретится, а значит $M'$ попросту никогда не остановится на данном входе. Итак, по определению полурешающей МТ построенная машина $M'$ является таковой для языка $\mathcal{L}$, а значит теорема доказана.
\end{proof} \leavevmode
\\
Построим по данному перечислимому языку $\mathcal{L}$ и его перечисляющей $k$-ленточной машине Тьюринга $M$ новую $(k + 3)$-ленточную перечисляющую $\mathcal{L}$ машину Тьюринга $M'$, которая каждое слово из языка выводит только один раз. \\
Пусть на лентах машины $M'$ с номерами от $1$ до $k$ эмулируется работа машины $M$, на $(k + 1)$ ленте хранится последнее выведенное машиной $M$ слово $w$, на $(k + 2)$ ленте, являющейся точным дубликатом выхода $M'$ (то есть, эта лента на каждой следующей итерации формирования вывода машины $M'$ должна в точности повторять содержимое $(k + 3)$-ей выходной ленты, но, если это требуется, в процессе проверки на наличие слова $w$ данную ленту можно изменять и записывать данные не только лишь один раз слева направо, что является неудобным формальным ограничением выходной ленты), ищется то самое слово $w$, и, если оно не найдено, выводится как на $(k + 2)$, так и на $(k + 3)$ ленту, а на последней $(k + 3)$ ленте печатается выход $M'$ (то есть эта лента --- выходная). Легко видеть, что $M'$ в самом деле перечисляет язык $\mathcal{L}$, но слова на её выходной ленте не повторяются, так как при формировании перечисления на $(k + 2)$-ой ленте происходит проверка наличия слова в перечислении, что воспрепятствует возникновению повторов. Таким образом, построение требуемой машины завершено.

\begin{lemma}
    Если существует МТ, перечисляющая язык $\mathcal{L}$ в лексикографическом порядке, то язык $\mathcal{L}$ разрешим.
\end{lemma}

\begin{proof}
    Пусть машина Тьюринга $M$ перечисляет $\mathcal{L}$ в лексикографическом порядке. Построим решающую машину $M'$ для $\mathcal{L}$ следующим образом: на входе $x$ машина $M'$ эмулирует работу $M$ с пустого входа и сравнивает выводимые $M$ слова со словом $x$ до тех пор, пока либо не произойдёт совпадение, либо $M$ не выведет слово длины больше, чем длина слова $x$. Если произошло совпадение, машина $M'$ останавливается на данном входе в состоянии $q_{\text{yes}}$, иначе, когда длина будет превышена, машина $M'$ останавливается в состоянии $q_{\text{no}}$. Из определения лексикографического порядка очевидно, что после перескока границы длины рассматриваемое слово более точно не встретится в перечислении, а значит, если оно не было найдено ранее, его нет в языке $\mathcal{L}$. Таким образом, построенная машина $M'$ удовлетворят определению решающей машины для $\mathcal{L}$, а значит язык $\mathcal{L}$ разрешим и лемма доказана.
\end{proof}

\section{Перечислимые языки}

\begin{theorem}
    Если язык и его дополнение перечислимы, то язык разрешим.
\end{theorem}

\begin{proof}
    Пусть язык $\mathcal{L}$ над алфавитом $\Sigma$ таков, что $\mathcal{L}$ и $\overline{\mathcal{L}}$ --- перечислимы. Тогда существуют две машины Тьюринга $M$ и $\overline{M}$, перечисляющие $\mathcal{L}$ и $\overline{\mathcal{L}}$ соответственно. Из определения перечисляющей МТ сразу следует, что любое слово $x$ над $\Sigma$ за конечное время появится либо на выходной ленте $M$, либо на выходной ленте $\overline{M}$, ведь $x$ лежит либо в $\mathcal{L}$, либо в $\overline{\mathcal{L}}$. Это означает, что за конечное время мы всегда сможем понять, запустив обе рассматриваемые машины, лежит слово в языке $\mathcal{L}$ или нет, что эквивалентно существованию решающей МТ для $\mathcal{L}$. Таким образом, язык $\mathcal{L}$ разрешим, что и требовалось.
\end{proof}

\begin{consequence}
    Перечислимые языки не замкнуты относительно дополнения.
\end{consequence}

\begin{proof}
    В самом деле, вновь рассмотрим язык \textbf{stop}. Ранее было доказано, что \textbf{stop} является полуразрешимым, но не является разрешимым. При этом полуразрешимость эквивалентна перечислимости, а значит язык \textbf{stop} является перечислимым, но не является разрешимым. Тогда его дополнение не может являться перечислимым в силу теоремы выше, а значит следствие показано.
\end{proof}

\begin{theorem}
    Множество перечислимо тогда и только тогда, когда оно есть область определения некоторой вычислимой функции.
\end{theorem}

\begin{proof}
    $(\Longleftarrow):$ Пусть функция $f$ является вычислимой. Без ограничения общности можно считать, что $f \colon S \to \mathbb{N}$, $S \subseteq \mathbb{N}$, так как каждому элементу области определения и области значений $f$, если это потребуется, мы можем взаимно-однозначно сопоставить натуральное число в бинарной записи. Из определения вычислимой функции вытекает, что существует машина Тьюринга $M_f$, вычисляющая $f$. Пусть машина $M_f'$ такова, что она останавливается в состоянии $q_{\text{yes}}$ на входе $x$, если $M_f$ останавливается на входе $x$ (то есть $x \in S$), и не останавливается на входе $x$, если $M_f$ не останавливается на входе $x$ (то есть в силу наших договорённостей это означает, что $x \not \in S$). Таким образом, для множества $S$ построена полурешающая машина Тьюринга $M_f'$, а значит $S$ перечислимо, что и требовалось. \\
    $(\Longrightarrow):$ Пусть множество $S$ перечислимо. Рассмотрим полухарактеристическую функцию $f \colon S \to \{1\}$ множества $S$, то есть $f(x) = 1$ тогда и только тогда, когда $x \in S$, а при всех остальных $x$ функция $f$ не определена. Увидим, что рассматриваемая функция вычислима. В самом деле, так как множество $S$ перечислимо, то для него существует полурешающая машина Тьюринга $M_S$. Тогда машина Тьюринга $M_S'$, которая в случае остановки $M_S$ в состоянии $q_{\text{yes}}$ печатает $1$, а в случае безостановочной работы $M_S$ также не останавливается, очевидным образом является вычисляющей МТ для $f$. Таким образом, множество $S$ является областью определения вычислимой функции $f$, а значит теорема доказана.
\end{proof}

\begin{theorem}
    Множество перечислимо тогда и только тогда, когда оно есть область значений некоторой вычислимой функции.
\end{theorem}

\begin{proof}
    $(\Longleftarrow):$ Пусть функция $f \colon S \to \mathbb{N}$, $S \subseteq \mathbb{N}$, является вычислимой. Тогда существует машина Тьюринга $M_f$, её вычисляющая. Построим машину Тьюринга $M$, перечисляющую область значений $f$. Для этого вновь создадим бесконечную таблицу, в первом столбце которой перечислены все натуральные аргументы, а в первой строке --- количества шагов, которые (или не более, чем которые) машина $M_f$ должна совершить при вычислении $f$ от соответствующего аргумента.
    \begin{table}[htbp]
    \centering
    \renewcommand{\arraystretch}{2}
    \begin{tabular}{|c|c|c|c|c|c|}
        \hline
        \diagbox{Аргумент}{Кол-во шагов} & 1000 & 2000 & 3000 & 4000 & $\dots$ \\
        \hline
        $1$ & \rdot{n11} & \rdot{n12} & \rdot{n13} & \rdot{n14} & \\
        \hline
        $2$ & \rdot{n21} & \rdot{n22} & \rdot{n23} & \rdot{n24} & \\
        \hline
        $3$ & \rdot{n31} & \rdot{n32} & \rdot{n33} & \rdot{n34} & \\
        \hline
        $4$ & \rdot{n41} & \rdot{n42} & \rdot{n43} & \rdot{n44} & \\
        \hline
        $\dots$ & & & & & \\
        \hline
    \end{tabular}
    \label{tab:snake_time_2}
    \end{table}

    \begin{tikzpicture}[remember picture, overlay, >=stealth]
    \tikzset{snake path/.style={red, thick, ->}}

    \draw[snake path] (n11) -- (n21);
    \draw[snake path] (n21) -- (n22);
    \draw[snake path] (n22) -- (n12);

    \draw[snake path] (n12) -- (n13);
    \draw[snake path] (n13) -- (n23);
    \draw[snake path] (n23) -- (n33);
    \draw[snake path] (n33) -- (n32);
    \draw[snake path] (n32) -- (n31);

    \draw[snake path] (n31) -- (n41);
    \draw[snake path] (n41) -- (n42);
    \draw[snake path] (n42) -- (n43);
    \draw[snake path] (n43) -- (n44);
    \draw[snake path] (n44) -- (n34);
    \draw[snake path] (n34) -- (n24);
    \draw[snake path] (n24) -- (n14);

    \draw[red, thick, dashed, ->] (n14) -- ++(1, 0); 
    \end{tikzpicture} \leavevmode
    \\
    Снова проходим эту таблицу змейкой так, как показано на рисунке (все клетки, как ранее было показано, будут посещены). Если в некоторой клетке для данного аргумента машина Тьюринга $M_f$ остановилась за требуемое число шагов, то машина $M$ печатает на ленте соответствующее вычисленное значение $f$, иначе ничего не печатает. Таким образом, для любого значения из области значений $f$ существует аргумент, на котором оно принимается, а значит за конечное число шагов вычисляющая МТ $M_f$ остановится на этом входе. Следовательно, в таблице существует клетка, в которую за конечное число шагов попадёт змейка, и $M$ напечатает рассматриваемое значение на ленте. Значение же не из области значений $f$ машина $M$, как легко видеть из построения, никогда не напечатает. Итак, для множества значений $f$ построена перечисляющая его МТ $M$, а значит оно перечислимо, что и требовалось. \\
    $(\Longrightarrow):$ Пусть множество $S$ перечислимо. Тогда существует машина Тьюринга $M_S$, его перечисляющая. Рассмотрим функцию $f$ такую, что $\forall i \in \mathbb{N}$ значение $f(i)$ есть $i$-ое слово на выходной ленте перечисляющей МТ $M_S$. Увидим, что эта функция вычислима. Действительно, машина Тьюринга $M$, которая на входе $i$ эмулирует работу $M_S$ на пустом входе и останавливается на $i$-ом слове, печатая его на ленту, является вычисляющей МТ для функции $f$. При этом ясно, что множество значений $f$ по определению перечисляющей МТ в точности совпадает с множеством $S$, а значит теорема доказана.
\end{proof}

\end{document}